\section{Related Work}
\label{sec:related-work}

\subsection{Lean and mathlib}

The formalization is implemented in Lean 4
\cite{demoura2015lean,demoura2021lean4,leanprover,tpil4} and builds on mathlib
\cite{mathlib2020,mathlibdocs,vandoorn2020maintaining,baanen2025growing}.  The project follows the mathlib-compatible
strategy of using existing topology, analysis, manifolds, finite-dimensional
calculus, and exterior-algebra infrastructure where possible.  The local box
Stokes layer is organized around mathlib-style divergence and differentiability
interfaces rather than a separate analytic foundation.

The theorem also exposes places where future mathlib APIs would be valuable.
In particular, the represented-to-native bridge would naturally depend on a
stable API for integrating differential forms on oriented manifolds with
boundary, together with chart-independence and boundary-orientation comparison
theorems.

\subsection{Prior Lean Stokes artifacts}

The repository treats \code{d0d1/lean-stokes-theorem} as prior work and a Lake
dependency reference \cite{leanstokes}.  That project is useful prior art for
Stokes formalization, especially around cubical or smooth-singular-cube
approaches.  Because of licensing considerations, this project uses external
developments as references rather than as copied source.  The present
contribution is different in emphasis: it formalizes a compact-support
coordinate-represented local-to-global route through chart boxes,
support-controlled partitions, half-space boundary boxes, zero-localized
representatives, and generated represented endpoints.

\subsection{Classical manifold treatments}

Classical treatments of manifolds and differential forms provide the
mathematical background for chartwise local arguments, partitions of unity,
compact support, boundary orientations, and Stokes' theorem
\cite{spivak1965calculus,munkres1991analysis,guillemin1974differential,bott1982differential,hirsch1976differential,warner1983foundations,abraham1988manifolds,lang1995differential,tu2011introduction,lee2013smooth,morita2001geometry}.
The formalization can be read as a detailed elaboration of the proof steps
that such texts compress.  The point is not that those steps are
mathematically surprising to a geometer.  The point is that their exact
dependencies become visible when the proof is checked by Lean.

This is especially visible for boundary charts.  Classical accounts freely
move between the local half-space model, the true boundary face, and the
auxiliary faces of a coordinate box.  In Lean these become separate theorems:
the lower face must be identified with the boundary contribution, the sign
must match the outward-first convention, and every artificial face must be
removed by a support theorem.

\subsection{Formalized mathematics and proof engineering}

The work also sits in a broader theorem-proving ecosystem.  Coq/Rocq,
Isabelle/HOL, and HOL Light have long provided reference points for formalized
mathematics and proof engineering
\cite{bertot2004coqart,nipkow2002isabelle,harrison2009handbook,wiedijk2006seventeen,harrison2014history}.
Within that ecosystem, formalized analysis and geometry have repeatedly
required substantial library and interface work, for example filters and
analysis in Isabelle/HOL, smooth manifolds in Isabelle/HOL, formalized Haar
measure, and the change-of-variables theorem in mathlib
\cite{holzl2013filters,immler2019smooth,vandoorn2021haar,gouezel2022change}.

Large formalizations often spend much of their effort on interfaces: finding
the right theorem boundary, replacing manual certificates by constructors, and
organizing local lemmas so that global statements have natural hypotheses.
Prominent formalization projects such as CompCert
\cite{leroy2009formal}, the formal four-color theorem
\cite{gonthier2008fourcolor}, the machine-checked odd-order theorem
\cite{gonthier2013machine}, the formal proof of the Kepler conjecture
\cite{hales2017kepler}, Schemes in Lean \cite{buzzard2022schemes},
formalized Galois theory \cite{browning2022galois}, and the Liquid Tensor
Experiment \cite{scholze2022liquid} illustrate that the
architecture of a large proof artifact is part of the contribution.  This
project is an example in differential geometry and analysis.  The intermediate
routes were not abandoned failures; they were stages in identifying which
fields were true mathematical hypotheses and which were internal proof
objects.  The first-principles theorem is the compressed result of that
interface work.

This distinction matters for evaluating formalized mathematics.  A
certificate-consuming theorem may be enough to close a proof internally, but it
does not yet present the theorem in the form a mathematician expects to use.
The present artifact treats assumption compression as part of the mathematical
work.  Covers, refinements, support facts, local regularity records, and
endpoint reconstruction are not left as user obligations once they can be
generated from compact support and chartwise smoothness.  This agrees with a
broader view of formalization as both theorem proving and mathematical
interface design \cite{avigad2018mechanization,baanen2022instance}.

The support and zero-localization layer is also an example of a common
phenomenon in formalization: a statement that is harmless on paper becomes
false or ill-shaped when all functions are total.  The solution is not merely a
technical patch.  It identifies the correct formal representative of the
informal mathematical object.
